\section{Discussion}
\label{sec:discussion}

\begin{table}[t]
\centering\scriptsize
\setlength{\tabcolsep}{3pt}
\renewcommand{\arraystretch}{1.0}
\caption{\textbf{Credit levels compared by where their likelihood comes from.}}
\resizebox{\columnwidth}{!}{%
\begin{tabular}{@{}llccccc@{}}
\toprule
\textbf{Level} & \textbf{Likelihood from} & \textbf{Env.} & \textbf{No oracle} & \textbf{Belief} & \textbf{Act gap} & \textbf{Maturity} \\
\midrule
Output & humans, verifiers & \hbnone & \hbhalf & \hbnone & \hbnone & \hbfull \\
Trajectory & verifier, PRM & \hbhalf & \hbhalf & \hbnone & \hbnone & \hbfull \\
Interaction (actions) & task reward & \hbfull & \hbfull & \hbnone & \hbnone & \hbfull \\
Interaction (memory ops) & task reward, teacher & \hbfull & \hbhalf & \hbhalf & \hbnone & \hbhalf \\
State (oracle or teacher) & closed world, teacher, target & \hbnone & \hbnone & \hbfull & \hbfull & \hbhalf \\
State (own observation) & returned observation & \hbfull & \hbfull & \hbfull & \hbhalf & \hbhalf \\
\bottomrule
\end{tabular}}
\tabnote{Env., the likelihood comes from the environment; No oracle, no external source of truth is needed; Belief, the written belief receives the gradient; Act gap, belief and action error are separated; Maturity, \hbfull\ many systems, \hbhalf\ a few, \hbnone\ none.}
\label{tab:levels}
\end{table}


Table~\ref{tab:levels} compares the credit levels by where their likelihood comes from; three observations follow.

\paragraph{The likelihood is the scarce term.}
Table~\ref{tab:learningmain} shows the state level taking its likelihood from a source the environment does not supply in every case but two: ReBel settles a predicate belief against the next observation, ABBEL reconstructs the previous one. Both beliefs are binary, both results are isolated, and the negative result beside them shows a dense environment-signed reward is unstable under group normalisation. That row is no longer empty; it is thin.

\paragraph{Auditability and expressiveness pull apart.}
Filters are validated in closed form but cannot introduce a factor they were not given, written beliefs can name a new factor but are rarely testable, and world models are testable but mute; hybrids that write the belief and check it parametrically are the convergence point \citep{song2026agent,zhang2026selfevolvingx,kumar2026towards}.

\paragraph{The belief--action gap is a training problem.}
Wherever belief and action are scored on the same episodes, prompting and oracles improve the belief while action--belief consistency barely moves \citep{deb2026stocktake,gao2026auditing}; the system that moves both trains belief and policy jointly \citep{singh2026agent}. Closing the gap requires the state-level and interaction-level signals of Table~\ref{tab:levels} in one objective, not a better belief alone.
